% Source: https://artofproblemsolving.com/community/c3h2431528_is_there_a_better_way_to_do_this_problem
% Stable identity: aops-topic-2431528-q01
The hexagon is a simple (non-self-intersecting) polygon, so the shoelace formula applies. With $(x_i,y_i)$ the vertices in order, the terms $x_iy_{i+1}-x_{i+1}y_i$ are
\begin{align*}
0\cdot(-4)-3\cdot0&=0, & 3\cdot(-3)-8\cdot(-4)&=23, & 8\cdot1-5\cdot(-3)&=23,\\
5\cdot6-6\cdot1&=24, & 6\cdot5-1\cdot6&=24, & 1\cdot0-0\cdot5&=0,
\end{align*}
with sum $94$. The area is $\dfrac{94}{2}=\boxed{47}$.

\begin{asy}
usepackage("amsmath");
size(5cm);
defaultpen(fontsize(10pt));
pair[] P={(0,0),(3,-4),(8,-3),(5,1),(6,6),(1,5)};
fill(P[0]--P[1]--P[2]--P[3]--P[4]--P[5]--cycle,rgb(0.85,0.9,1));
for(int i=-1;i<=9;++i) draw((i,-5)--(i,7),gray(0.85));
for(int j=-5;j<=7;++j) draw((-1,j)--(9,j),gray(0.85));
draw(P[0]--P[1]--P[2]--P[3]--P[4]--P[5]--cycle,linewidth(1));
string[] L={"$(0,0)$","$(3,-4)$","$(8,-3)$","$(5,1)$","$(6,6)$","$(1,5)$"};
pair[] D={W,S,E,W,NE,NW};
for(int i=0;i<6;++i){dot(P[i]); label(L[i],P[i],D[i]);}
\end{asy}

(Check with Pick's theorem: the hexagon has $B=6$ lattice points on its boundary, one at each vertex, since every side has coprime coordinate differences, and $I=45$ interior lattice points, so its area is $I+\dfrac{B}{2}-1=45+3-1=47$.)
